\section{Method}\label{sec:method}
Let $f\colon [0,1]^2 \to \mathbb{R}_{\ge 0}$ be the heat map and let
$x_1,\dots,x_n$ be the probe locations. After $k$ probes we fit the estimate
\begin{equation}\label{eq:estimate}
  \hat f_k(x) = \sum_{i=1}^{k} w_i(x)\, f(x_i),
  \qquad \sum_{i=1}^{k} w_i(x) = 1,
\end{equation}
and choose the next probe where the local variation of $\hat f_k$ is largest.

\subsection{Error bound}\label{sec:bound}
Suppose $f$ is $L$-Lipschitz and supported on $s$ disks of radius $r$.
Then the greedy rule reaches error $\varepsilon$ after
$O\!\left(s \log(1/\varepsilon)\right)$ probes, against
$O\!\left(\varepsilon^{-2}\right)$ for uniform sampling~\cite{donoho2006}.

The proof follows the covering argument of~\cite{candes2006}: each hot spot is
found after a bounded number of probes, and the remaining probes refine the
estimate (\ref{eq:estimate}) inside the disks.
